\section{Unsmoothed moments and the covariance of actual returns}
\label{sec:unsmoothed-moments}
The covariance of a weak Gaussian limit need not be the limit of the
normalized second moments. We establish that identification for the
actual return record, including a bounded-variation insertion at any
return. The argument also gives fourth moments for the unsmoothed
collision sums and improves the two-sided stopping estimate. It uses
only the collision splitting, not decay of correlations for $F_R^*$.

Write $\|u\|_{\mathcal V}=\|u\|_\infty+\|u\|_{\mathrm{BV}}$.
All expectations without a subscript in this section are with respect
to the collision probability $\nu$. Constants are uniform in $R\in I$.
The smoothing operators preserve means and have the bounds in
Lemma~\ref{lem:bv-smoothing}.

\subsection{Decoupling at a deterministic collision gap}
\begin{lemma}[A finite-product decoupling estimate]
\label{lem:bv-finite-product}
For $2\le q\le4$, integers $t_1\le\cdots\le t_q$, and
$1\le p<q$, bounded-variation functions $u_1,\ldots,u_q$ satisfy
\begin{align}
 &\left|\mathbb E\prod_{j=1}^q u_j\circ T_R^{t_j}
 -\left(\mathbb E\prod_{j=1}^p u_j\circ T_R^{t_j}\right)
  \left(\mathbb E\prod_{j=p+1}^q u_j\circ T_R^{t_j}\right)\right|
 \notag\\
 &\hspace{18mm}\le C\exp\{-\gamma(t_{p+1}-t_p)\}
                         \prod_{j=1}^q\|u_j\|_{\mathcal V}.
 \label{eq:bv-finite-product}
\end{align}
The times may be negative or repeated. The constants $C,\gamma>0$
can be chosen for all the indicated values of $q$ and $p$.
\end{lemma}
\begin{proof}
By multilinearity assume $\|u_j\|_{\mathcal V}\le1$; a zero norm
makes the assertion immediate. Replace every $u_j$ by
$u_{j,\epsilon}=\mathsf S_\epsilon u_j$. Telescoping the products,
invariance of $\nu$, and the uniform supremum bounds show that the
change in the left-hand expression is at most $C_q\epsilon$,
independently of all the times. In particular, no variation norm of
$u_j\circ T_R^{t_j}$ is estimated.

Translate $t_1$ to zero. Put $d_j=t_{j+1}-t_j$ and use the transfer
convention of Lemma~\ref{lem:collision-spectral-input}. The smoothed
joint expectation is
\[
 \ell\bigl(M_{u_{q,\epsilon}}\mathcal L_R^{d_{q-1}}
 M_{u_{q-1,\epsilon}}\cdots
 \mathcal L_R^{d_1}M_{u_{1,\epsilon}}\nu\bigr).
\]
Replace the power across the chosen gap by $\Pi_R=\nu\ell$.
The resulting expression is exactly the product of the two smoothed
expectations in \eqref{eq:bv-finite-product}. The difference is bounded
by $C_q\epsilon^{-2q}\vartheta^{d_p}$. Indeed each smooth multiplier
costs at most $C\epsilon^{-2}$, every other power is uniformly bounded,
and the complementary power across the gap costs $C\vartheta^{d_p}$.
For a gap of length zero use
$\mathcal L_R^0-\Pi_R=I-\Pi_R$, with its uniform bound.

Choose $\epsilon=\tfrac14\vartheta^{d_p/(2q+1)}$. Smoothing and
operator errors are at most $C_q\vartheta^{d_p/(2q+1)}$. The finite
parameter cover permits $\gamma=-\log\vartheta/9$ after changing $C$.
Invariance justifies the initial translation also for negative times.
\end{proof}

\begin{proposition}[Fourth moments without smoothing]
\label{prop:unsmoothed-fourth}
Let $u$ be real, centered, and $\|u\|_{\mathcal V}\le B$.
For every deterministic integer interval of length $m\ge1$,
\begin{equation}\label{eq:unsmoothed-fourth}
 \mathbb E\left|\sum_{j=r}^{r+m-1}u\circ T_R^j\right|^4\le C_Bm^2,
 \qquad
 \mathbb E\max_{0\le l\le m}
 \left|\sum_{j=r}^{r+l-1}u\circ T_R^j\right|^4\le C_Bm^2.
\end{equation}
Both estimates hold in the reverse order of the interval and under
any nonnegative initial density bounded by a fixed constant. They hold
componentwise, and hence in Euclidean norm, for $h_R$.
\end{proposition}
\begin{proof}
Put $c_u(a)=\mathbb E[u(u\circ T_R^a)]$. For $a,b,c\ge0$, apply
Lemma~\ref{lem:bv-finite-product} to four copies of $u$ at times
$0,a,a+b,a+b+c$. Splitting across the middle gap bounds
\[
 \left|\mathbb E[u(u\circ T_R^a)(u\circ T_R^{a+b})
                   (u\circ T_R^{a+b+c})]-c_u(a)c_u(c)\right|
       \le C_B e^{-\gamma b}.
\]
Splitting off the first or last centered factor bounds the fourfold
expectation by $C_Be^{-\gamma a}$ and $C_Be^{-\gamma c}$,
respectively. Also $|c_u(a)c_u(c)|\le C_Be^{-\gamma(a+c)}$.
Taking the minimum of these three bounds gives
\begin{align}\label{eq:four-product-gap}
 &\left|\mathbb E[u(u\circ T_R^a)(u\circ T_R^{a+b})
                   (u\circ T_R^{a+b+c})]-c_u(a)c_u(c)\right|
 \le C_B e^{-\gamma\max\{a,b,c\}}.
\end{align}
This includes repeated times.

Expand the fourth power and order its four indices; their multiplicity
is at most $24$. For each initial index the sum of
$e^{-\gamma(a+c)}$ over the two outer gaps is bounded, while the middle
gap has at most $m$ choices. This contributes $O(m^2)$. The sum over
all three nonnegative gaps of $e^{-\gamma\max\{a,b,c\}}$ is finite,
so the remaining contribution is $O(m)$. This proves the first bound.

For the second enlarge $m$ to a power of two $M<2m$. Every prefix is
a disjoint union of at most one dyadic interval of each length.
Consequently its absolute sum is bounded by the sum, over scales, of
the maximum absolute interval sum at that scale. At interval length
$l$, the fourth power of this maximum has expectation at most
$C_B(M/l)l^2=C_BMl$. Minkowski's inequality therefore gives an $L^4$
norm at most
\[
 C_B\sum_{j=0}^{\log_2 M}(M^2 2^{-j})^{1/4}
       \le C_B M^{1/2}.
\]
There is no logarithmic loss. Stationarity translates deterministic
intervals; reversal gives suffixes of the same interval. Domination
by a bounded initial density proves the remaining assertions.
\end{proof}

\subsection{Fourth-moment windows and two-sided stopping}
Retain $N^\pm_{j,R}$, $m_j=\lfloor j/c_*\rfloor$ and
$\mathsf H_{r,s,R}$ from Section~\ref{sec:marked-return-band}. Define
\begin{equation}\label{eq:moment-stopped-pair}
 Z_{n,k,R}=\mathsf H_{N^-_{k,R},N^+_{n-k,R},R},\qquad
 W_{n,k,R}=\mathsf H_{m_k,m_{n-k},R}.
\end{equation}
Thus $Z_{n,k,R}(y)=J_{n,R}((F_R^*)^{-k}y)-n\bar G_R$ almost
everywhere. Zero-length backward and forward intervals are empty.

\begin{lemma}[Fourth-moment return-clock window]
\label{lem:fourth-clock-window}
For $j\ge0$, $b\ge2$, and either sign,
\begin{equation}\label{eq:fourth-clock-window}
 \nu_R^*\{|N^\pm_{j,R}-m_j|>b\}
       \le C\frac{(j+b)^2}{b^4}.
\end{equation}
\end{lemma}
\begin{proof}
The case $j=0$ is immediate. The centered visit count
$\sum_{i=1}^L(\eta_R-c_*)\circ T_R^{\pm i}$ has fourth moment at
most $CL^2$ under $\nu_R^*$, by
Proposition~\ref{prop:unsmoothed-fourth}. The exact visit equivalences
in Lemma~\ref{lem:two-sided-clock-window} turn the upper and lower
return-clock events into deviations of these counts. At integer
thresholds within a bounded distance of $m_j\pm b$, the required
deviation is at least $c_*b-O(1)$ and $L\le C(j+b)$. Markov's
inequality proves \eqref{eq:fourth-clock-window}. A negative lower
threshold contributes no event. Increase $C$ for bounded $b$.
\end{proof}

\begin{theorem}[Moment control of the actual marked stopping]
\label{thm:marked-L2-stopping}
Uniformly for $n\ge2$, $0\le k\le n$, and $R\in I$,
\begin{align}
 \int (|Z_{n,k,R}|^4+|W_{n,k,R}|^4)\dd\nu_R^*&\le Cn^2,
       \label{eq:actual-marked-fourth}\\
 \left\|\frac{Z_{n,k,R}-W_{n,k,R}}{\sqrt n}
                  \right\|_{L^2(\nu_R^*)}&\le Cn^{-1/6}.
       \label{eq:actual-marked-L2}
\end{align}
For every real $v\in\R^4$ and insertion $a$ of
\eqref{eq:marked-norms}, the stopping comparison also satisfies
\begin{equation}\label{eq:improved-marked-stopping}
 \left|C^{[k],a}_{n,R}(v)
       -\int_M a e^{iv\cdot W_{n,k,R}/\sqrt n}\dd\nu\right|
       \le C M_a(1+|v|)n^{-2/9}.
\end{equation}
\end{theorem}
\begin{proof}
The deterministic interval has length at most $Cn$, so its fourth
moment follows from Proposition~\ref{prop:unsmoothed-fourth} and
$\nu_R^*\le c_*^{-1}\nu$. For the random interval put
$Q=N^-_{k,R}+N^+_{n-k,R}$. Invariance of the actual return map gives
$N^-_{j,R}\overset{\rm law}=N^+_{j,R}$ under $\nu_R^*$: use
$N^-_{j,R}((F_R^*)^jx)=N^+_{j,R}(x)$. Hence
Theorem~\ref{thm:cumulative-return-tail} and a union bound imply
\[
 \nu_R^*(Q>L)\le C e^{an-cL/2}.
\]
No independence of the two clocks is used. Fix $A>2a/c+2$.
On $Q\le An$, $|Z_{n,k,R}|$ is bounded by the sum of the backward
and forward collision maxima up to $\lceil An\rceil$. Their fourth
moments are $O(n^2)$. On $Q>An$, boundedness of $h_R$ gives
$|Z_{n,k,R}|\le C Q$. Integrating the displayed exponential tail
bounds that fourth-moment contribution by
$C(1+n)^4e^{-c' n}$ for some $c'>0$. This proves
\eqref{eq:actual-marked-fourth}.

Let $2\le b\le n$ and let $\mathcal E_b$ be the union of the two
clock-window exceptions. Lemma~\ref{lem:fourth-clock-window} gives
$\nu_R^*(\mathcal E_b)\le Cn^2/b^4$. Off this event, the difference
$Z_{n,k,R}-W_{n,k,R}$ is bounded by maxima on two deterministic
collision windows of length $O(b)$, including windows that cross zero.
Their second moments are $O(b)$ by the fourth-moment maximal estimate.
On the exceptional event, Cauchy--Schwarz and
\eqref{eq:actual-marked-fourth} give
\[
 \int |Z_{n,k,R}-W_{n,k,R}|^2\one_{\mathcal E_b}\dd\nu_R^*
       \le Cn\,\nu_R^*(\mathcal E_b)^{1/2}\le Cn^2/b^2.
\]
Thus the normalized squared $L^2$ error is at most
$C(b/n+n/b^2)$. Taking $b=\lceil n^{2/3}\rceil$ proves
\eqref{eq:actual-marked-L2}.

For the characteristic comparison, the exceptional event costs only
$CM_a n^2/b^4$. On its complement, the same deterministic windows and
$|e^{ix}-e^{iy}|\le|x-y|$ give $CM_a|v|\sqrt{b/n}$.
Choose instead $b=\lceil n^{5/9}\rceil$. Both powers are
$n^{-2/9}$. The recentering identity
\eqref{eq:marked-recentering} and $|a|\nu\le c_*M_a\nu_R^*$
complete the proof. All rounding exceptions are absorbed in $C$.
\end{proof}

\subsection{Quadratic insertions and the induced covariance}
The following estimate supplies the moment identification that is not
asserted by the characteristic-function proof in
Remark~\ref{rem:gaussian-versus-raw}.

\begin{lemma}[A centered insertion in a quadratic collision sum]
\label{lem:quadratic-collision-insertion}
Let $a$ be a complex bounded-variation function on $M$, and let
$\alpha=\int a\dd\nu$. For all $r,s\ge0$, with $m=r+s$,
\begin{align}
 \left|\int a\mathsf H_{r,s,R}\dd\nu\right|
       &\le C(\|a\|_\infty+\|a\|_{\mathrm{BV}}),
       \label{eq:linear-insertion-moment}\\
 \left\|\int a\mathsf H_{r,s,R}\otimes\mathsf H_{r,s,R}\dd\nu
                          -\alpha m\Gamma_R\right\|
       &\le C(\|a\|_\infty+\|a\|_{\mathrm{BV}}).
       \label{eq:quadratic-insertion-moment}
\end{align}
The constants do not depend on the location of the origin in the interval.
\end{lemma}
\begin{proof}
Center $a$ by subtracting $\alpha$. Each component of $h_R$ is
centered. The two-factor version of
Lemma~\ref{lem:bv-finite-product} bounds
$\int(a-\alpha)h_R\circ T_R^i\dd\nu$ by
$C\|a\|_{\mathcal V}e^{-\gamma|i|}$. Summation proves the first
assertion.

For components $u,v$ of $h_R$, sort the three times $0,i,j$ and
denote their two consecutive gaps by $d_1,d_2$. All three factors
$a-\alpha,u,v$ are centered. Splitting off the first and the last
factor in \eqref{eq:bv-finite-product} gives
\[
 \left|\int(a-\alpha)(u\circ T_R^i)(v\circ T_R^j)\dd\nu\right|
       \le C\|a\|_{\mathcal V}e^{-\gamma\max\{d_1,d_2\}}.
\]
Because $\max(d_1,d_2)\ge\max(|i|,|j|)/2$, the right side is
summable over all $(i,j)\in\Z^2$. The centered insertion therefore
changes the quadratic moment of any such interval by at most
$C\|a\|_{\mathcal V}$. Finally the exponentially summable collision
covariances give
\[
 \left\|\int\mathsf H_{r,s,R}\otimes\mathsf H_{r,s,R}\dd\nu
                                  -m\Gamma_R\right\|\le C.
\]
This is the finite covariance expansion with the lag weights $m-l$;
the sum of $l$ times the absolute lag covariances is finite. Multiplying
by $\alpha$ proves the second assertion. Real and imaginary parts
handle complex $a$ without any positivity assumption.
\end{proof}

\begin{theorem}[Moments at an arbitrary actual return mark]
\label{thm:marked-quadratic-moments}
Put $U_{n,R}=J_{n,R}-n\bar G_R$. For $a$ as in
\eqref{eq:marked-norms}, uniformly for $n\ge2$ and $0\le k\le n$,
\begin{align}
 \left|\int_{Y_R^*}a((F_R^*)^kx)\frac{U_{n,R}(x)}{\sqrt n}\dd\nu(x)\right|
 &\le C\left(M_an^{-1/6}+(M_a+V_a)n^{-1/2}\right),
       \label{eq:marked-first-moment}\\
 \left\|\int_{Y_R^*}a((F_R^*)^kx)
       \frac{U_{n,R}(x)\otimes U_{n,R}(x)}n\dd\nu(x)-\alpha_a D_R\right\|
 &\le C\left(M_an^{-1/6}+(M_a+V_a)n^{-1}\right).
       \label{eq:marked-second-moment}
\end{align}
These are moments of the finite complex measure of
\eqref{eq:marked-measure}, using $\nu|_{Y_R^*}$, not its normalization.
In particular
\begin{equation}\label{eq:actual-covariance-rate}
 \sup_{R\in I}\left\|\frac1n\int U_{n,R}\otimes U_{n,R}\dd\nu_R^*
                                               -D_R\right\|
       \le C n^{-1/6}.
\end{equation}
\end{theorem}
\begin{proof}
Recenter at $y=(F_R^*)^kx$. The weighted moments then involve
$a(y)Z_{n,k,R}(y)$ and $a(y)Z_{n,k,R}(y)\otimes Z_{n,k,R}(y)$.
Theorem~\ref{thm:marked-L2-stopping} changes the normalized first
moment by at most $CM_an^{-1/6}$. For the second use
\[
 \|z\otimes z-w\otimes w\|\le |z-w|(|z|+|w|)
\]
and Cauchy--Schwarz. Equations~\eqref{eq:actual-marked-fourth} and
\eqref{eq:actual-marked-L2} show that its normalized change is also
at most $CM_an^{-1/6}$. The factor $c_*$ in passing from
$\nu|_{Y_R^*}$ to $\nu_R^*$ is fixed.

Apply Lemma~\ref{lem:quadratic-collision-insertion} to
$W_{n,k,R}=\mathsf H_{m_k,m_{n-k},R}$. Its length satisfies
$m_k+m_{n-k}=n/c_*+O(1)$ uniformly in $k$. Since
$D_R=\Gamma_R/c_*$, division by $\sqrt n$ or $n$ gives
\eqref{eq:marked-first-moment} and \eqref{eq:marked-second-moment}.
For \eqref{eq:actual-covariance-rate} take $a=s_R=c_*^{-1}\eta_R$,
whose mass is one and whose two norms are uniformly bounded. Invariance
of $F_R^*$ and the Kac mean give $\int U_{n,R}\dd\nu_R^*=0$.
Thus the last display is convergence of the actual covariance matrices.
\end{proof}

\begin{corollary}[The induced Green--Kubo formula with Cesaro weights]
\label{cor:induced-cesaro-covariance}
Let $g_R=G_R-\bar G_R$ and
$C^*_{j,R}=\int g_R\otimes(g_R\circ(F_R^*)^j)\dd\nu_R^*$.
Then
\begin{equation}\label{eq:induced-cesaro-covariance}
 \left\|C^*_{0,R}+\sum_{j=1}^{n-1}\left(1-\frac jn\right)
          (C^*_{j,R}+(C^*_{j,R})^{\mathsf T})-D_R\right\|
       \le C n^{-1/6}.
\end{equation}
\end{corollary}
\begin{proof}
The exponential one-return moment makes each term finite. Expand the
second moment of $\sum_{j=0}^{n-1}g_R\circ(F_R^*)^j$ and use stationarity
under $F_R^*$. The left matrix before subtracting $D_R$ is exactly
$n^{-1}\int U_{n,R}\otimes U_{n,R}\dd\nu_R^*$. Apply
\eqref{eq:actual-covariance-rate}.
\end{proof}

\begin{corollary}[Moments under the unchanged marked-state event]
\label{cor:marked-conditional-moments}
Use the event $A_{n,k,R}$ and variation $V_{n,R}$ of
Corollary~\ref{cor:marked-state-conditioning}, and put
$p_{n,R}=\nu_R^*(E_{n,R})>0$. The conditional first moment of
$U_{n,R}/\sqrt n$ is bounded by
\[
 \frac C{p_{n,R}}\left(n^{-1/6}+(1+V_{n,R})n^{-1/2}\right).
\]
Its conditional second moment about the original center $n\bar G_R$
differs from $D_R$ by at most
\[
 \frac C{p_{n,R}}\left(n^{-1/6}+(1+V_{n,R})n^{-1}\right).
\]
Consequently, if $p_{n,R}\ge c n^{-\beta}$ and
$V_{n,R}\le Cn^\kappa$, the latter difference tends to zero when
$\beta<1/6$ and $\beta+\kappa<1$. The conditional mean tends to zero,
and the conditional covariance tends to $D_R$, under the stronger
conditions $\beta<1/6$ and $\beta+\kappa<1/2$.
\end{corollary}
\begin{proof}
Take $a=c_*^{-1}\one_{E_{n,R}}$ in
Theorem~\ref{thm:marked-quadratic-moments}. Its mass is $p_{n,R}$.
Divide the two estimates by this unchanged probability. Subtracting
the tensor square of the conditional mean identifies the conditional
covariance. The strict exponent inequalities prove convergence.
\end{proof}

\paragraph{Role in raw local inversion.}
Theorems~\ref{thm:marked-L2-stopping} and
\ref{thm:marked-quadratic-moments} concern the unsmoothed physical
record, not a replacement process. They identify its covariance by
actual second moments, establish uniform integrability of its squared
normalized size, and strengthen the stopping part of Theorem C.
In particular all the rare marked-state events in
Corollary~\ref{cor:marked-state-conditioning} also satisfy the new
conditional moment conclusions. Cesaro convergence in
\eqref{eq:induced-cesaro-covariance} does not assert absolute summability
of the induced correlations.

Positive definiteness is a distinct assertion. The estimates above give
no positive lower bound for the smallest eigenvalue of $D_R$ and no
periodic-evaluable representative of its zero-variance coboundary.
Theorem~\ref{thm:joint-uniform-nondegeneracy} supplies positive
definiteness by a separate measurable phase argument, not from these
moment bounds. Likewise \eqref{eq:improved-marked-stopping} is not a complementary
frequency estimate. The physical central radius remains
$2n^{-99/200}$, or $2n^{1/200}$ after rescaling. The annulus beyond it,
the complete raw critical and singular branch sum, and the weighted
tails for an exact physical observation retain their roles in
Theorem~\ref{thm:LLT}. No conditioning event has been replaced.
